\section{Categorical representations and outer automorphisms}

Recall the category
\[
    \operatorname{Rep}(G,\mathcal{C}),
\]
whose objects are pairs \((X,\rho)\), where \(X\in \mathcal{C}\) and
\[
    \rho:G\longrightarrow \operatorname{Aut}_{\mathcal{C}}(X)
\]
is a group homomorphism. A morphism
\[
    f:(X,\rho)\longrightarrow (X',\rho')
\]
is a morphism \(f:X\to X'\) in \(\mathcal{C}\) such that, for every \(g\in G\),
\[
    \rho'(g)\circ f=f\circ \rho(g).
\]

\begin{lemma}
    A morphism
    \[
        f:(X,\rho)\longrightarrow (X',\rho')
    \]
    in \(\operatorname{Rep}(G,\mathcal{C})\) is an isomorphism if and only if the
    underlying morphism
    \[
        f:X\longrightarrow X'
    \]
    is an isomorphism in \(\mathcal{C}\).
\end{lemma}

\begin{proof}
    If \(f\) is an isomorphism in \(\operatorname{Rep}(G,\mathcal{C})\), then its
    underlying morphism in \(\mathcal{C}\) is clearly an isomorphism.

    Conversely, suppose that \(f:X\to X'\) is an isomorphism in \(\mathcal{C}\).
    Since \(f\) is a morphism of representations, for every \(g\in G\) we have
    \[
        \rho'(g)\circ f=f\circ \rho(g).
    \]
    Multiplying by \(f^{-1}\), we obtain
    \[
        f^{-1}\circ \rho'(g)=\rho(g)\circ f^{-1}.
    \]
    Hence \(f^{-1}:X'\to X\) is also a morphism of representations. Therefore
    \(f\) is an isomorphism in \(\operatorname{Rep}(G,\mathcal{C})\).
\end{proof}

\begin{definition}
    Let \((X,\rho)\) and \((X,\rho')\) be two representations of \(G\) on the same
    object \(X\). We say that \((X,\rho)\) and \((X,\rho')\) are isomorphic if there
    exists an automorphism
    \[
        f:X\longrightarrow X
    \]
    in \(\mathcal{C}\) such that, for every \(g\in G\),
    \[
        \rho'(g)=f\circ \rho(g)\circ f^{-1}.
    \]
    In this case, we also say that \(\rho\) and \(\rho'\) are uniformly conjugate.
\end{definition}

If two representations \((X,\rho)\) and \((X',\rho')\) are isomorphic, we write
    \[
        (X,\rho)\cong (X',\rho').
    \]

Let \((X,\rho)\) and \((Y,\sigma)\) be two \(\mathcal{C}\)-representations of \(G\).
Then the group
\[
    \operatorname{Aut}_{\mathcal{C}}(X)\times \operatorname{Aut}_{\mathcal{C}}(Y)
\]
acts on the set
\[
    \operatorname{Hom}_{\mathcal{C}}(X,Y)
\]
by
\[
    (\alpha,\beta)\cdot f
    :=
    \beta\circ f\circ \alpha^{-1},
\]
where
\[
    \alpha\in \operatorname{Aut}_{\mathcal{C}}(X),\qquad
    \beta\in \operatorname{Aut}_{\mathcal{C}}(Y),\qquad
    f\in \operatorname{Hom}_{\mathcal{C}}(X,Y).
\]

The diagonal group homomorphism
\[
    G\longrightarrow
    \operatorname{Aut}_{\mathcal{C}}(X)\times \operatorname{Aut}_{\mathcal{C}}(Y)
\]
is given by
\[
    g\longmapsto \bigl(\rho(g),\sigma(g)\bigr).
\]
Therefore it induces an action of \(G\) on \(\operatorname{Hom}_{\mathcal{C}}(X,Y)\):
\[
    g\cdot f
    :=
    \sigma(g)\circ f\circ \rho(g)^{-1},
    \qquad
    g\in G,\quad f\in \operatorname{Hom}_{\mathcal{C}}(X,Y).
\]

Equivalently, the following diagram commutes:
\[
\begin{tikzcd}
    X \arrow[r,"f"] \arrow[d,"\rho(g)"']
        & Y \arrow[d,"\sigma(g)"] \\
    X \arrow[r,"g\cdot f"']
        & Y
\end{tikzcd}
\]
where
\[
    g\cdot f=\sigma(g)\circ f\circ \rho(g)^{-1}.
\]

\begin{remark}
    A morphism \(f:X\to Y\) is a morphism of representations
    \[
        f:(X,\rho)\longrightarrow (Y,\sigma)
    \]
    if and only if it is fixed by the above \(G\)-action, that is,
    \[
        g\cdot f=f
        \qquad
        \text{for all } g\in G.
    \]
    Indeed, this condition is equivalent to
    \[
        \sigma(g)\circ f=f\circ \rho(g),
    \]
    which is exactly the equivariance condition.
\end{remark}

\begin{lemma}
    Let \((X,\rho)\) and \((Y,\sigma)\) be two objects of
    \(\operatorname{Rep}(G,\mathcal{C})\). Then the set of morphisms
    \[
        (X,\rho)\longrightarrow (Y,\sigma)
    \]
    in \(\operatorname{Rep}(G,\mathcal{C})\) is precisely the set of fixed points
    of \(G\) in its action on
    \[
        \operatorname{Hom}_{\mathcal{C}}(X,Y).
    \]
\end{lemma}

\begin{proof}
    Recall that \(G\) acts on \(\operatorname{Hom}_{\mathcal{C}}(X,Y)\) by
    \[
        g\cdot f
        :=
        \sigma(g)\circ f\circ \rho(g)^{-1}.
    \]
    Thus \(f\) is fixed by every \(g\in G\) if and only if
    \[
        \sigma(g)\circ f\circ \rho(g)^{-1}=f
    \]
    for every \(g\in G\). This is equivalent to
    \[
        \sigma(g)\circ f=f\circ \rho(g),
    \]
    which is exactly the condition that \(f\) is a morphism of representations
    \[
        f:(X,\rho)\longrightarrow (Y,\sigma).
    \]
    Therefore the morphisms in \(\operatorname{Rep}(G,\mathcal{C})\) are precisely
    the \(G\)-fixed points of \(\operatorname{Hom}_{\mathcal{C}}(X,Y)\).
\end{proof}

Assume further that \(\mathcal{C}\) is a category whose objects are sets and whose
morphisms are maps between sets.

Let \((X,\rho)\) be a \(\mathcal{C}\)-representation of \(G\). We write
\[
    g\cdot x := \rho(g)(x),
    \qquad x\in X,\quad g\in G.
\]
Thus the representation \(\rho:G\to \operatorname{Aut}_{\mathcal{C}}(X)\) gives an
action of \(G\) on the set \(X\).

A morphism
\[
    f:(X,\rho)\longrightarrow (Y,\sigma)
\]
is a morphism \(f:X\to Y\) in \(\mathcal{C}\) satisfying
\[
    f(g\cdot x)=g\cdot f(x)
\]
for all \(x\in X\) and \(g\in G\). Such a map is also called a \(G\)-equivariant
map.

\begin{definition}
    Let \(G\) act on a set \(X\). For \(x\in X\), the stabilizer of \(x\) is the
    subgroup
    \[
        G_x:=\{\,g\in G\mid g\cdot x=x\,\}.
    \]
\end{definition}

\begin{lemma}
    Let \((X,\rho)\) be a \(\mathcal{C}\)-representation of \(G\). Then:
    \begin{enumerate}
        \item For every \(g\in G\) and every \(x\in X\),
        \[
            G_{g\cdot x}=gG_xg^{-1}.
        \]

        \item If
        \[
            f:(X,\rho)\longrightarrow (Y,\sigma)
        \]
        is an isomorphism of \(G\)-representations, then
        \[
            G_{f(x)}=G_x
        \]
        for every \(x\in X\).
    \end{enumerate}
\end{lemma}

\begin{proof}
    \textup{(1)} Let \(h\in G\). Then
    \[
        h\in G_{g\cdot x}
    \]
    if and only if
    \[
        h\cdot (g\cdot x)=g\cdot x.
    \]
    This is equivalent to
    \[
        (g^{-1}hg)\cdot x=x,
    \]
    that is,
    \[
        g^{-1}hg\in G_x.
    \]
    Therefore
    \[
        h\in gG_xg^{-1}.
    \]
    Hence
    \[
        G_{g\cdot x}=gG_xg^{-1}.
    \]

    \textup{(2)} Since \(f\) is an isomorphism of \(G\)-representations, it is
    \(G\)-equivariant and bijective. For \(g\in G\), we have
    \[
        g\in G_x
        \iff g\cdot x=x.
    \]
    Applying \(f\) and using equivariance gives
    \[
        f(g\cdot x)=f(x)
        \iff g\cdot f(x)=f(x).
    \]
    Hence
    \[
        g\in G_x
        \iff g\in G_{f(x)}.
    \]
    Therefore
    \[
        G_x=G_{f(x)}.
    \]
\end{proof}

Let \(G\) be a group. Denote by
\[
    \operatorname{Aut}(G)
\]
the automorphism group of \(G\).

For each \(g\in G\), define a map
\[
    \operatorname{Inn}(g):G\longrightarrow G
\]
by
\[
    \operatorname{Inn}(g)(x)=gxg^{-1},
    \qquad x\in G.
\]
This is called the inner automorphism of \(G\) induced by \(g\).

\begin{definition}
    The group of inner automorphisms of \(G\) is
    \[
        \operatorname{Inn}(G)
        :=
        \{\,\operatorname{Inn}(g)\mid g\in G\,\}.
    \]
\end{definition}

\begin{remark}
    For every \(g\in G\), the map \(\operatorname{Inn}(g)\) is an automorphism of
    \(G\). Indeed,
    \[
        \operatorname{Inn}(g)(xy)
        =
        gxyg^{-1}
        =
        (gxg^{-1})(gyg^{-1})
        =
        \operatorname{Inn}(g)(x)\operatorname{Inn}(g)(y),
    \]
    and its inverse is
    \[
        \operatorname{Inn}(g^{-1}).
    \]
\end{remark}

\begin{proposition}
    There is a group isomorphism
    \[
        \operatorname{Inn}(G)\cong G/Z(G),
    \]
    where
    \[
        Z(G):=\{\,z\in G\mid zg=gz \text{ for all } g\in G\,\}
    \]
    is the center of \(G\).
\end{proposition}

\begin{proof}
    Define
    \[
        \Phi:G\longrightarrow \operatorname{Aut}(G)
    \]
    by
    \[
        \Phi(g)=\operatorname{Inn}(g).
    \]
    We first check that \(\Phi\) is a group homomorphism. For \(g,h\in G\) and
    \(x\in G\), we have
    \[
        \bigl(\operatorname{Inn}(g)\circ \operatorname{Inn}(h)\bigr)(x)
        =
        \operatorname{Inn}(g)(hxh^{-1})
        =
        ghxh^{-1}g^{-1}
        =
        (gh)x(gh)^{-1}.
    \]
    Hence
    \[
        \operatorname{Inn}(g)\circ \operatorname{Inn}(h)
        =
        \operatorname{Inn}(gh),
    \]
    so
    \[
        \Phi(gh)=\Phi(g)\Phi(h).
    \]
    Thus \(\Phi\) is a group homomorphism.

    By definition,
    \[
        \operatorname{Im}(\Phi)=\operatorname{Inn}(G).
    \]
    Now we compute the kernel. We have
    \[
        g\in \operatorname{Ker}(\Phi)
    \]
    if and only if
    \[
        \operatorname{Inn}(g)=\operatorname{id}_G.
    \]
    This means that for every \(x\in G\),
    \[
        gxg^{-1}=x.
    \]
    Equivalently,
    \[
        gx=xg
    \]
    for every \(x\in G\). Hence
    \[
        \operatorname{Ker}(\Phi)=Z(G).
    \]
    Therefore, by the first isomorphism theorem,
    \[
        G/Z(G)\cong \operatorname{Im}(\Phi)=\operatorname{Inn}(G).
    \]
\end{proof}

\begin{proposition}
    The group of inner automorphisms is a normal subgroup of the automorphism
    group:
    \[
        \operatorname{Inn}(G)\trianglelefteq \operatorname{Aut}(G).
    \]
    More precisely, for every \(g\in G\) and every \(a\in \operatorname{Aut}(G)\),
    one has
    \[
        a\circ \operatorname{Inn}(g)\circ a^{-1}
        =
        \operatorname{Inn}(a(g)).
    \]
\end{proposition}

\begin{proof}
    Let \(a\in \operatorname{Aut}(G)\), \(g\in G\), and \(x\in G\). Then
    \[
        \bigl(a\circ \operatorname{Inn}(g)\circ a^{-1}\bigr)(x)
        =
        a\Bigl(g\,a^{-1}(x)\,g^{-1}\Bigr).
    \]
    Since \(a\) is a group homomorphism, this equals
    \[
        a(g)\,a(a^{-1}(x))\,a(g^{-1}).
    \]
    Hence
    \[
        \bigl(a\circ \operatorname{Inn}(g)\circ a^{-1}\bigr)(x)
        =
        a(g)\,x\,a(g)^{-1}.
    \]
    Therefore
    \[
        a\circ \operatorname{Inn}(g)\circ a^{-1}
        =
        \operatorname{Inn}(a(g)).
    \]
    Since \(a(g)\in G\), we have
    \[
        \operatorname{Inn}(a(g))\in \operatorname{Inn}(G).
    \]
    Thus every conjugate of an inner automorphism by an automorphism is again an
    inner automorphism. Hence
    \[
        \operatorname{Inn}(G)\trianglelefteq \operatorname{Aut}(G).
    \]
\end{proof}

Let \(G\) be a group. Recall that
\[
    \operatorname{Aut}(G)
\]
denotes the automorphism group of \(G\).

Let \((X,\rho)\) be a \(\mathcal{C}\)-representation of \(G\), where
\[
    \rho:G\longrightarrow \operatorname{Aut}_{\mathcal{C}}(X)
\]
is a group homomorphism. For \(a\in \operatorname{Aut}(G)\), define a new
representation
\[
    {}^{a}(X,\rho):=(X,{}^{a}\rho),
\]
where
\[
    {}^{a}\rho:G\longrightarrow \operatorname{Aut}_{\mathcal{C}}(X)
\]
is given by
\[
    {}^{a}\rho=\rho\circ a^{-1}.
\]
Equivalently, for \(g\in G\),
\[
    {}^{a}\rho(g)=\rho(a^{-1}(g)).
\]

\begin{lemma}
    For every \(a\in \operatorname{Aut}(G)\), the pair
    \[
        {}^{a}(X,\rho)=(X,{}^{a}\rho)
    \]
    is again a \(\mathcal{C}\)-representation of \(G\).
\end{lemma}

\begin{proof}
    Since \(a^{-1}:G\to G\) is a group automorphism and
    \(\rho:G\to \operatorname{Aut}_{\mathcal{C}}(X)\) is a group homomorphism,
    their composite
    \[
        {}^{a}\rho=\rho\circ a^{-1}
    \]
    is again a group homomorphism. Therefore
    \[
        {}^{a}\rho:G\longrightarrow \operatorname{Aut}_{\mathcal{C}}(X)
    \]
    defines a \(\mathcal{C}\)-representation of \(G\).
\end{proof}

\begin{lemma}
    Let
    \[
        f:(X,\rho)\longrightarrow (X',\rho')
    \]
    be a morphism of \(\mathcal{C}\)-representations of \(G\). Then, for every
    \(a\in \operatorname{Aut}(G)\), the same underlying morphism
    \[
        f:X\longrightarrow X'
    \]
    induces a morphism of representations
    \[
        f:{}^{a}(X,\rho)\longrightarrow {}^{a}(X',\rho').
    \]
\end{lemma}

\begin{proof}
    Since \(f:(X,\rho)\to (X',\rho')\) is a morphism of representations, for every
    \(h\in G\) one has
    \[
        \rho'(h)\circ f=f\circ \rho(h).
    \]
    Now let \(g\in G\). Applying the above identity to \(h=a^{-1}(g)\), we get
    \[
        \rho'(a^{-1}(g))\circ f
        =
        f\circ \rho(a^{-1}(g)).
    \]
    By definition of the twisted representations, this is precisely
    \[
        {}^{a}\rho'(g)\circ f
        =
        f\circ {}^{a}\rho(g).
    \]
    Hence \(f\) is a morphism
    \[
        {}^{a}(X,\rho)\longrightarrow {}^{a}(X',\rho').
    \]
\end{proof}

Therefore \(\operatorname{Aut}(G)\) acts on the set of isomorphism classes of
\(\mathcal{C}\)-representations of \(G\) by
\[
    a\cdot [(X,\rho)]
    :=
    [\,{}^{a}(X,\rho)\,].
\]

\begin{remark}
    This is indeed a left action. For \(a,b\in \operatorname{Aut}(G)\), we have
    \[
        {}^{a}({}^{b}\rho)
        =
        (\rho\circ b^{-1})\circ a^{-1}
        =
        \rho\circ b^{-1}a^{-1}
        =
        \rho\circ (ab)^{-1}
        =
        {}^{ab}\rho.
    \]
    Hence
    \[
        {}^{a}\bigl({}^{b}(X,\rho)\bigr)
        =
        {}^{ab}(X,\rho).
    \]
\end{remark}

For \(g\in G\), recall that the inner automorphism induced by \(g\) is
\[
    \operatorname{Inn}(g):G\longrightarrow G,
    \qquad
    x\longmapsto gxg^{-1}.
\]
Thus
\[
    \operatorname{Inn}(g)^{-1}=\operatorname{Inn}(g^{-1}).
\]

Let \((X,\rho)\) be a \(\mathcal{C}\)-representation of \(G\). Then
\[
    {}^{\operatorname{Inn}(g)}\rho(h)
    =
    \rho\bigl(\operatorname{Inn}(g)^{-1}(h)\bigr)
    =
    \rho(g^{-1}hg)
\]
for every \(h\in G\).

\begin{proposition}
    For every \(g\in G\), the morphism
    \[
        \rho(g):X\longrightarrow X
    \]
    is an isomorphism of representations
    \[
        \rho(g):
        {}^{\operatorname{Inn}(g)}(X,\rho)
        \xrightarrow{\ \cong\ }
        (X,\rho).
    \]
\end{proposition}

\begin{proof}
    Since \(\rho(g)\in \operatorname{Aut}_{\mathcal{C}}(X)\), the underlying
    morphism \(\rho(g):X\to X\) is an isomorphism in \(\mathcal{C}\).

    It remains to check equivariance. For \(h\in G\), we need to prove
    \[
        \rho(h)\circ \rho(g)
        =
        \rho(g)\circ {}^{\operatorname{Inn}(g)}\rho(h).
    \]
    The left-hand side is
    \[
        \rho(h)\circ \rho(g)=\rho(hg).
    \]
    On the other hand,
    \[
        \rho(g)\circ {}^{\operatorname{Inn}(g)}\rho(h)
        =
        \rho(g)\circ \rho(g^{-1}hg)
        =
        \rho(hg).
    \]
    Hence
    \[
        \rho(h)\circ \rho(g)
        =
        \rho(g)\circ {}^{\operatorname{Inn}(g)}\rho(h).
    \]
    Therefore \(\rho(g)\) is an isomorphism of representations
    \[
        {}^{\operatorname{Inn}(g)}(X,\rho)
        \cong
        (X,\rho).
    \]
\end{proof}

\begin{corollary}
    Every inner automorphism of \(G\) acts trivially on the set of isomorphism
    classes of \(\mathcal{C}\)-representations of \(G\).
\end{corollary}

\begin{proof}
    By the previous proposition, for every \(g\in G\) and every representation
    \((X,\rho)\), we have an isomorphism
    \[
        {}^{\operatorname{Inn}(g)}(X,\rho)
        \cong
        (X,\rho).
    \]
    Hence
    \[
        \operatorname{Inn}(g)\cdot [(X,\rho)]
        =
        [(X,\rho)].
    \]
    Therefore inner automorphisms act trivially on isomorphism classes.
\end{proof}

Since
\[
    \operatorname{Inn}(G)\trianglelefteq \operatorname{Aut}(G),
\]
the action of \(\operatorname{Aut}(G)\) on isomorphism classes descends to an
action of the outer automorphism group
\[
    \operatorname{Out}(G)
    :=
    \operatorname{Aut}(G)/\operatorname{Inn}(G).
\]

Thus
\[
    \operatorname{Out}(G)
\]
acts naturally on the set of isomorphism classes of \(\mathcal{C}\)-representations
of \(G\).
